\cleardoublepage
\thispagestyle{empty}
\begin{rpgdark}[PROTOCOLO DE MESA]
\centering
{\textgoth{\fontsize{24}{26}\selectfont PROTOCOLO}}\\[0.25em]
{\scshape\bfseries\small QUATRO MODOS DE ENTRAR NO ARQUIVO}\\[0.35em]
{\color{rpggold}\rule{0.55\textwidth}{0.7pt}}
\end{rpgdark}

\vspace{0.5em}
\noindent{\small\itshape Este é um aparato de leitura, não um jogo completo. A obra pode ser interrompida a qualquer momento; a rota não muda o texto, só decide por onde você entra.}

\vspace{0.6em}
\begin{rpgbox}[1 · LEITURA LINEAR]
{\color{rpgblood}\faDiceD20}\enspace \textbf{Siga a ordem das páginas.} Do Sertão (1914) ao Epílogo. É o caminho recomendado para a primeira leitura: acompanha Joaquim e a propagação da entidade em ordem temporal.
\end{rpgbox}

\begin{rpgbox}[2 · ROTA INVESTIGATIVA]
{\color{rpgblood}\faDiceD20}\enspace \textbf{Comece pela Parte 4 (Investigação Lúcia, 2026).} Descubra as anomalias junto com a psiquiatra forense e volte às Partes 1, 2 e 3 para remontar as provas como uma autópsia.
\end{rpgbox}

\begin{rpgbox}[3 · NAVEGAÇÃO HIPERTEXTUAL]
{\color{rpgblood}\faDiceD20}\enspace \textbf{Siga as rotas ao final de cada fragmento.} Cada salto reflete a fragmentação do arquivo. Nenhum caminho é um beco: de qualquer fragmento existe percurso até o Epílogo.
\end{rpgbox}

\begin{rpgbox}[4 · ROTAS CURADAS]
{\color{rpgblood}\faDiceD20}\enspace \textbf{Sequências pré-selecionadas.} Rota do Terror: CONT-03 $\to$ CONT-04 $\to$ MEM-12 $\to$ CONT-05. Rota da Compaixão: MEM-01 $\to$ MEM-05 $\to$ MEM-10 $\to$ MEM-15. Rota Forense: DOC-01 $\to$ LUC-01 $\to$ DOC-16 $\to$ CONT-05.
\end{rpgbox}

\vspace{0.5em}
\begin{rpgbox}[ROTA ALEATÓRIA · role 1d20]
\centering
\small
\begin{tabular}{@{}cl@{}}
\toprule
\textbf{d20} & \textbf{Entrada} \\
\midrule
1--5 & Sertão \\
6--10 & Colônia \\
11--15 & Diário de Oliveira e Laudos \\
16--20 & Investigação Lúcia \\
\bottomrule
\end{tabular}

\vspace{0.4em}
{\itshape A Parte 5 (Contaminação) deve ser lida por último em qualquer percurso. Ela não é sobre a história. Ela é sobre você.}
\end{rpgbox}

\vspace{0.5em}
\begin{rpgbox}[CÓDIGOS DO ARQUIVO]
\small
\textcolor{mem}{\textbf{MEM}} Memória · \textcolor{doc}{\textbf{DOC}} Documento · \textcolor{luc}{\textbf{LUC}} Lucidez · \textcolor{cont}{\textbf{CONT}} Contaminação
\end{rpgbox}
\clearpage
